This section specifies the APIs provided and required by the major components identified in the Component View. The interfaces are categorized into \textbf{Hardware Abstraction Interfaces} (internal to the Hybrid Client) and \textbf{Network Interfaces} (between Client and Cloud).

\subsection{Hardware Abstraction Interfaces}
These interfaces are exposed by the \textbf{Native Bridge (Capacitor)} and consumed by the \textbf{Sensor Controller} component within the web runtime. They allow the TypeScript application logic to access device sensors.

\paragraph{Interface: IGeolocationProvider}
\begin{itemize}
    \item \textbf{Description:} Provides access to the device's GPS hardware.
    \item \textbf{Provided By:} Native Bridge (\texttt{@capacitor/geolocation})
    \item \textbf{Used By:} Sensor Controller
\end{itemize}

\begin{table}[H]
\centering
\renewcommand{\arraystretch}{1.4}
\begin{tabular}{|p{0.22\textwidth}|p{0.25\textwidth}|p{0.18\textwidth}|p{0.25\textwidth}|}
\hline
\textbf{Method Signature} & \textbf{Input Parameters} & \textbf{Return Value} & \textbf{Description} \\
\hline
\texttt{getCurrentPosition()} & \texttt{options: GeoOptions} & \texttt{Promise<Position>} & Returns the single current coordinate (Lat, Long, Altitude). \\
\hline
\texttt{watchPosition()} & \texttt{options: GeoOptions, callback: (pos) => void} & \texttt{CallbackID} & Starts a continuous background location tracking session. The callback is invoked every time the OS receives a location update. \\
\hline
\texttt{clearWatch()} & \texttt{id: CallbackID} & \texttt{void} & Stops the tracking session to save battery. \\
\hline
\end{tabular}
\caption{IGeolocationProvider Interface}
\end{table}

\paragraph{Interface: IMotionSensor}
\begin{itemize}
    \item \textbf{Description:} Provides high-frequency access to the accelerometer for pothole detection.
    \item \textbf{Provided By:} Native Bridge (\texttt{@capacitor/motion})
    \item \textbf{Used By:} Sensor Controller (Inertial Analyzer)
\end{itemize}

\begin{table}[H]
\centering
\renewcommand{\arraystretch}{1.4}
\begin{tabular}{|p{0.22\textwidth}|p{0.25\textwidth}|p{0.18\textwidth}|p{0.25\textwidth}|}
\hline
\textbf{Method Signature} & \textbf{Input Parameters} & \textbf{Return Value} & \textbf{Description} \\
\hline
\texttt{addListener()} & \texttt{eventName: 'accel', callback: (event) => void} & \texttt{PluginListenerHandle} & Subscribes to the raw accelerometer stream. The event contains $x, y, z$ acceleration vectors used to detect vertical shocks. \\
\hline
\texttt{removeAllListeners()} & \texttt{void} & \texttt{void} & Unsubscribes from the sensor stream. \\
\hline
\end{tabular}
\caption{IMotionSensor Interface}
\end{table}

\subsection{Cloud Network Interfaces (BaaS API)}
These interfaces represent the communication contract between the \textbf{Hybrid Client} and the \textbf{Serverless Cloud}. Since the architecture uses a BaaS, these are primarily RESTful endpoints automatically generated by Supabase from the database schema, plus RPC (Remote Procedure Calls) for Edge Functions.

\paragraph{Interface: IAuthService}
\begin{itemize}
    \item \textbf{Description:} Manages user sessions.
    \item \textbf{Protocol:} HTTPS / REST
    \item \textbf{Provided By:} Auth Service (Supabase GoTrue)
\end{itemize}

\begin{table}[H]
\centering
\renewcommand{\arraystretch}{1.4}
\begin{tabular}{|p{0.25\textwidth}|p{0.25\textwidth}|p{0.20\textwidth}|p{0.20\textwidth}|}
\hline
\textbf{Endpoint / Method} & \textbf{Input Payload} & \textbf{Response} & \textbf{Description} \\
\hline
\texttt{POST /auth/v1/signup} & \texttt{\{ email, password \}} & \texttt{\{ user, session, jwt \}} & Registers a new user and returns the access token. \\
\hline
\texttt{POST /auth/v1/token} & \texttt{\{ email, password, grant\_type \}} & \texttt{\{ access\_token, refresh\_token \}} & Authenticates a user (Login). \\
\hline
\texttt{POST /auth/v1/logout} & Header: \texttt{Bearer <token>} & \texttt{204 No Content} & Invalidates the current session. \\
\hline
\end{tabular}
\caption{IAuthService Interface}
\end{table}

\paragraph{Interface: ISpatialDataStore}
\begin{itemize}
    \item \textbf{Description:} Provides CRUD access to geospatial entities. Access is filtered by RLS policies.
    \item \textbf{Protocol:} HTTPS / REST
    \item \textbf{Provided By:} Spatial Database (PostgREST)
\end{itemize}

\begin{table}[H]
\centering
\renewcommand{\arraystretch}{1.4}
\begin{tabular}{|p{0.25\textwidth}|p{0.25\textwidth}|p{0.20\textwidth}|p{0.20\textwidth}|}
\hline
\textbf{Endpoint / Method} & \textbf{Input Payload} & \textbf{Response} & \textbf{Description} \\
\hline
\texttt{GET /rest/v1/paths} & Query Params: \texttt{select=*, geometry=...} & \texttt{JSON Array} & Retrieves public paths. Supports spatial filtering (e.g., retrieving paths within a map viewport). \\
\hline
\texttt{POST /rest/v1/activities} & \texttt{\{ user\_id, geometry\_linestring, ... \}} & \texttt{201 Created} & Uploads a recorded activity trace. This action triggers the \texttt{func\_path\_merging} Edge Function asynchronously. \\
\hline
\texttt{POST /rest/v1/feedbacks} & \texttt{\{ path\_id, rating, comment \}} & \texttt{201 Created} & Submits a user report for a specific path. \\
\hline
\end{tabular}
\caption{ISpatialDataStore Interface}
\end{table}

\paragraph{Interface: ISearchEngine (RPC)}
\begin{itemize}
    \item \textbf{Description:} Exposes the complex path ranking logic as a remote function call.
    \item \textbf{Protocol:} HTTPS / RPC
    \item \textbf{Provided By:} Edge Functions
\end{itemize}

\begin{table}[H]
\centering
\renewcommand{\arraystretch}{1.4}
\begin{tabular}{|p{0.25\textwidth}|p{0.25\textwidth}|p{0.20\textwidth}|p{0.20\textwidth}|}
\hline
\textbf{Function Name} & \textbf{Input Parameters} & \textbf{Return Value} & \textbf{Description} \\
\hline
\texttt{rpc/search\_best\_paths} & \texttt{\{ origin\_lat, origin\_lon, dest\_lat, dest\_lon \}} & \texttt{JSON Array <Path>} & Invokes the routing algorithm. Returns a list of paths ordered by the dynamic \textbf{Suitability Score} (calculated based on distance and road quality). \\
\hline
\end{tabular}
\caption{ISearchEngine Interface}
\end{table}

\subsection{External Interfaces}

\paragraph{Interface: IWeatherProvider}
\begin{itemize}
    \item \textbf{Description:} External API used to enrich path data with meteorological conditions.
    \item \textbf{Provided By:} Third-party provider (e.g., OpenWeatherMap)
    \item \textbf{Used By:} \texttt{func\_weather\_enricher} (Edge Function)
\end{itemize}

\begin{table}[H]
\centering
\renewcommand{\arraystretch}{1.4}
\begin{tabular}{|p{0.25\textwidth}|p{0.15\textwidth}|p{0.20\textwidth}|p{0.30\textwidth}|}
\hline
\textbf{Endpoint} & \textbf{Method} & \textbf{Params} & \textbf{Description} \\
\hline
\texttt{/data/2.5/weather} & \texttt{GET} & \texttt{lat, lon, appid} & Fetches current weather data (Temperature, Rain volume, Wind) for a specific coordinate. \\
\hline
\end{tabular}
\caption{IWeatherProvider Interface}
\end{table}